% Section 13.6: The External Field Approximation.
% Source: physical PDF pp. 579--585 (printed pp. 556--562).
% Starts with the Section 13.6 heading; ends with the paragraph immediately before Problems.
% Equations: (13.6.1)--(13.6.16). Figures: 13.4--13.5.
% Footnotes: four ordinary numbered notes, including the section-title note.
% Status: source-reviewed against rendered pages; no unresolved uncertainties.

\subsection[The External Field Approximation]{The External Field Approximation\texorpdfstring{\protect\footnotemark}{}}
\footnotetext{This section lies somewhat out of the book's main line of
development, and may be omitted in a first reading.}

It is intuitively obvious that a heavy charged particle like the nucleus of
an atom acts approximately like the source of a classical external field. In
this section we will see how to justify this approximation, and will gain
some idea of its limitations.

Consider a Feynman diagram or a part of a Feynman diagram in which
a heavy charged particle passing through the diagram from the initial to
the final state emits $N$ off-shell photons with four-momenta $q_1,q_2,\ldots,q_N$
and polarization indices $\mu_1,\mu_2,\ldots,\mu_N$. The sum of all such graphs or
subgraphs (not including the $N$ photon propagators) yields an amplitude
\begin{align}
&\int d^4x_1\,d^4x_2\cdots d^4x_N\,
e^{-iq_1\cdot x_1}e^{-iq_2\cdot x_2}\cdots e^{-iq_N\cdot x_N}
\notag\\
&\qquad\cdot\bra{p',\sigma'}T\{J^{\mu_1}(x_1),J^{\mu_2}(x_2),
\ldots,J^{\mu_N}(x_N)\}\ket{p,\sigma}
\notag\\
&\qquad\equiv\mathcal G^{\mu_1\mu_2\cdots\mu_N}_{\sigma',\sigma}
(q_1,q_2,\ldots,q_N;p).
\tag{13.6.1}\label{eq:13.6.1}
\end{align}
with the matrix element calculated including all interactions in which
the heavy particle may participate, including strong nuclear forces. This
amplitude has a multiple pole at $q_1,q_2,\ldots,q_N\to0$, arising from terms in
the matrix elements of the product of currents in which the intermediate
states consist of just the same heavy particle as in the initial and final states.
This multiple pole dominates \eqref{eq:13.6.1} when all components of $q_1,q_2,\ldots,q_N$
are small compared with all energies and momenta associated with the
dynamics of the (perhaps composite) heavy particle. In this case the
methods of Section 10.2 give\footnote{In perturbation theory, the denominators come from the denominators of the propagators:
\[
(p'+q_1+\cdots+q_r)^2+m^2-i\epsilon
\to2p'\cdot(q_1+\cdots+q_r)-i\epsilon
\to2p\cdot(q_1+\cdots+q_r)-i\epsilon,
\]
while the numerators of the propagators provide factors $\sum_\sigma uu^\dagger$
that together with the photon emission vertex matrices yield the matrix
elements \eqref{eq:13.6.3}. The matrix $\mathcal G^\mu$ differs from the matrix
$G^\mu$ of the previous section by a factor $2p^0$.}
\begin{align}
\mathcal G^{\mu_1\mu_2\cdots\mu_N}_{\sigma',\sigma}
(q_1,q_2,\ldots,q_N;p)
&\to\frac{(-i)^{N-1}}{2p^0(2\pi)^3}(2\pi)^4
\notag\\
&\quad\cdot\delta^4(p'+q_1+q_2+\cdots+q_N-p)
\sum_{\sigma_1,\sigma_2,\ldots,\sigma_{N-1}}
\notag\\
&\quad\cdot
\frac{\mathcal G^{\mu_1}_{\sigma',\sigma_1}(p)
\mathcal G^{\mu_2}_{\sigma_1,\sigma_2}(p)\cdots
\mathcal G^{\mu_N}_{\sigma_{N-1},\sigma}(p)}
{[2p\cdot q_1-i\epsilon]\!\!
[2p\cdot(q_1+q_2)-i\epsilon]\!\cdots
[2p\cdot(q_1+\cdots+q_{N-1})-i\epsilon]}
\notag\\
&\quad+\text{permutations}.
\tag{13.6.2}\label{eq:13.6.2}
\end{align}
where
\begin{equation}
\frac{\mathcal G^\mu_{\sigma',\sigma}(p)}{2p^0(2\pi)^3}
\equiv\bra{p,\sigma'}J^\mu(0)\ket{p,\sigma}
\tag{13.6.3}\label{eq:13.6.3}
\end{equation}
and `$+\text{permutations}$' indicates that we are to sum over all permutations
of the $N$ photons. For applications to atomic systems it is important
to recognize that \eqref{eq:13.6.1} applies for particles of arbitrary spin that have
strong as well as electromagnetic interactions, like atomic nuclei.

We also note that for particles of arbitrary spin and charge $Ze$, the
matrix elements of the electric current between states of equal four-momenta are\footnote{This is most easily proved by first noting that in the Lorentz frame in which the particle is at
rest, rotational invariance requires that the matrix elements of the current have vanishing space
components and a time component proportional to $\delta_{\sigma',\sigma}$, with no other dependence on $\sigma$ or $\sigma'$.
The constant of proportionality is supplied by Eq.~\eqref{eq:10.6.3}, and a Lorentz transformation then
gives Eq.~\eqref{eq:13.6.4}.}
\begin{equation}
\bra{p,\sigma'}J^\mu(0)\ket{p,\sigma}
=\frac{Ze\,p^\mu\delta_{\sigma',\sigma}}{p^0(2\pi)^3},
\tag{13.6.4}\label{eq:13.6.4}
\end{equation}
so that
\begin{equation}
\mathcal G^\mu_{\sigma',\sigma}(p)=2Ze\,p^\mu\delta_{\sigma',\sigma}.
\tag{13.6.5}\label{eq:13.6.5}
\end{equation}
The important thing about Eq.~\eqref{eq:13.6.5} is that these matrices all commute,
so their product can be factored out of the sum over permutations:
\begin{align}
\mathcal G^{\mu_1\mu_2\cdots\mu_N}_{\sigma',\sigma}
(q_1,q_2,\ldots,q_N;p)
&\to
\frac{(-i)^{N-1}(Ze)^N p^{\mu_1}p^{\mu_2}\cdots p^{\mu_N}}
{p^0(2\pi)^3}(2\pi)^4
\notag\\
&\quad\cdot
\delta^4(p'+q_1+q_2+\cdots+q_N-p)\delta_{\sigma',\sigma}
\notag\\
&\quad\cdot\Biggl[
\frac{1}{[p\cdot q_1-i\epsilon]
[p\cdot(q_1+q_2)-i\epsilon]\cdots
[p\cdot(q_1+\cdots+q_{N-1})-i\epsilon]}
\notag\\
&\qquad+\text{permutations}\Biggr].
\tag{13.6.6}\label{eq:13.6.6}
\end{align}
To leading order in the $q$'s the delta function here may be written
\begin{align}
\delta^4(p'+q_1+\cdots+q_N-p)={}&p^0
\delta^3(\mathbf p'+\mathbf q_1+\cdots+\mathbf q_N-\mathbf p)
\notag\\
&\cdot\delta\bigl(p\cdot(q_1+\cdots+q_N)\bigr).
\tag{13.6.7}\label{eq:13.6.7}
\end{align}
Fortunately, it turns out that the result of summing over permutations
here is much simpler than the individual terms. For
$p\cdot(q_1+\cdots+q_N)=0$, we have
\begin{align}
&\left[
\frac{1}{[p\cdot q_1-i\epsilon]
[p\cdot(q_1+q_2)-i\epsilon]\cdots
[p\cdot(q_1+\cdots+q_{N-1})-i\epsilon]}
+\text{permutations}\right]
\notag\\
&\qquad=(2i\pi)^{N-1}\delta(p\cdot q_1)\delta(p\cdot q_2)
\cdots\delta(p\cdot q_{N-1}).
\tag{13.6.8}\label{eq:13.6.8}
\end{align}
For instance, for $N=2$ this reads:
\begin{equation*}
\frac{1}{[p\cdot q_1-i\epsilon]}+\frac{1}{[p\cdot q_2-i\epsilon]}
=\frac{1}{[p\cdot q_1-i\epsilon]}+\frac{1}{[-p\cdot q_1-i\epsilon]}
=2i\pi\delta(p\cdot q_1).
\end{equation*}
The general result \eqref{eq:13.6.8} can be obtained most easily as the Fourier
transform of the identity
\begin{equation*}
\theta(\tau_1-\tau_2)\theta(\tau_2-\tau_3)\cdots
\theta(\tau_{N-1}-\tau_N)+\text{permutations}=1.
\end{equation*}
Inserting Eq.~\eqref{eq:13.6.8} in Eq.~\eqref{eq:13.6.6} gives our final result for the amplitude
\eqref{eq:13.6.1}:
\begin{align}
\mathcal G^{\mu_1\mu_2\cdots\mu_N}_{\sigma',\sigma}
(q_1,q_2,\ldots,q_N;p)
&\to
(Ze)^N(2\pi)^N\delta_{\sigma',\sigma}
p^{\mu_1}p^{\mu_2}\cdots p^{\mu_N}
\notag\\
&\quad\cdot
\delta^3(\mathbf p'+\mathbf q_1+\mathbf q_2+\cdots+\mathbf q_N-\mathbf p)
\notag\\
&\quad\cdot
\delta(p\cdot q_1)\delta(p\cdot q_2)\cdots\delta(p\cdot q_N).
\tag{13.6.9}\label{eq:13.6.9}
\end{align}
This result applies to relativistic as well as slowly moving heavy particles,
and can be used to derive the `Weizs\"acker--Williams' approximation~\hyperref[ref:13.9]{[9]} for
charged-particle scattering. In the special case of a non-relativistic heavy
charged particle, with $|\mathbf p|\ll p^0$, Eq.~\eqref{eq:13.6.9} further simplifies to
\begin{align}
\mathcal G^{\mu_1\mu_2\cdots\mu_N}_{\sigma',\sigma}
(q_1,q_2,\ldots,q_N;p)
&\to
(Ze)^N(2\pi)^N n^{\mu_1}n^{\mu_2}\cdots n^{\mu_N}
\notag\\
&\quad\cdot
\delta^3(\mathbf p'+\mathbf q_1+\mathbf q_2+\cdots+\mathbf q_N-\mathbf p)
\notag\\
&\quad\cdot
\delta(q_1^0)\delta(q_2^0)\cdots\delta(q_N^0)\delta_{\sigma',\sigma},
\tag{13.6.10}\label{eq:13.6.10}
\end{align}
where $n$ is a unit time-like vector
\begin{equation*}
n^0=1,\qquad \mathbf n=0.
\end{equation*}

Now suppose that a single heavy non-relativistic particle of charge $Ze$
with normalized momentum space wave function $\chi_\sigma(\mathbf p)$ appears in both
the initial and final states. Using the Fourier representation of the delta
function in Eq.~\eqref{eq:13.6.9}, the matrix element of $\mathcal G$ in this state is
\begin{align}
&\int d^3p\,d^3p'\,\chi^*_{\sigma'}(\mathbf p')\chi_\sigma(\mathbf p)
\notag\\
&\qquad\cdot
\mathcal G^{\mu_1\mu_2\cdots\mu_N}_{\sigma',\sigma}
(q_1,q_2,\ldots,q_N;p)
\notag\\
&\qquad\to
\int d^3X\sum_\sigma|\psi_\sigma(\mathbf X)|^2
\prod_{r=1}^N 2\pi Ze\,n^{\mu_r}\delta(q_r^0)e^{-i\mathbf q_r\cdot\mathbf X}.
\tag{13.6.11}\label{eq:13.6.11}
\end{align}
where $\psi_\sigma(\mathbf X)$ is the coordinate space wave function:
\begin{equation}
\psi_\sigma(\mathbf X)=(2\pi)^{-3/2}\int d^3p\,
\chi_\sigma(\mathbf p)e^{i\mathbf p\cdot\mathbf X}.
\tag{13.6.12}\label{eq:13.6.12}
\end{equation}
Because of the factorization in Eq.~\eqref{eq:13.6.11}, the effect of including a
heavy charged particle in this state is then the same as that of adding
any numbers of a new kind of vertex in the momentum space Feynman
rules, in which light Dirac particles of charge $-e$ such as electrons interact
with an external field, with each such vertex contributing to the overall
amplitude a factor\footnote{The first factor here is the usual factor of $i$ accompanying the constants in the interaction
Lagrangian of the heavy charged particle in the Feynman rules.} (now including the photon propagator and electron--
photon vertex)
\begin{align}
i\int d^4q\left[\frac{-i}{(2\pi)^4}\frac{1}{q^2-i\epsilon}\right]
\left[2\pi Ze\,n_\mu\delta(q^0)e^{-i\mathbf q\cdot\mathbf X}\right]
\left[(2\pi)^4e\gamma^\mu\delta^4(k-k'-q)\right].
\tag{13.6.13}\label{eq:13.6.13}
\end{align}
where $k$ and $k'$ are the initial and final electron four-momenta. The
complete scattering amplitude must then be averaged over the heavy
particle position $\mathbf X$, with weight function $\sum_\sigma|\psi_\sigma(\mathbf X)|^2$. The factor \eqref{eq:13.6.13}
is the same as would be produced by a new term in the interaction
Lagrangian
\begin{equation}
\mathcal L_{\mathrm{ext}}(x)=\mathcal A_\mu(x)J_e^\mu(x),
\tag{13.6.14}\label{eq:13.6.14}
\end{equation}
where $J_e^\mu\equiv-ie\bar\Psi\gamma^\mu\Psi$ is the electric current of the electrons, and $\mathcal A^\mu$ is an
external vector potential
\begin{equation}
\mathcal A^\mu(x)=\frac{1}{(2\pi)^4}\int d^4q\,e^{iq\cdot x}
\left[\frac{2\pi Ze\,n^\mu\delta(q^0)e^{-i\mathbf q\cdot\mathbf X}}
{q^2-i\epsilon}\right].
\tag{13.6.15}\label{eq:13.6.15}
\end{equation}
This, of course, is just the usual Coulomb potential:
\begin{equation}
\mathcal A^0(x)=\frac{Ze}{4\pi|\mathbf x-\mathbf X|},
\qquad \boldsymbol{\mathcal A}(x)=0.
\tag{13.6.16}\label{eq:13.6.16}
\end{equation}
If there is more than one heavy charged particle (as in a molecule) we
must express $\mathcal A^\mu(x)$ as a sum of terms like \eqref{eq:13.6.16}, each with its own
charge $Ze$ and position $\mathbf X$.

It is useful to keep in mind what diagrams we are summing in using the
external field approximation. Consider the interaction of a single electron
(relativistic or non-relativistic) with a single heavy charged particle such
as a proton or deuteron. If we ignore all other interactions, then the
Feynman diagrams for the scattering of the electron due to its interaction
with the external field are just those with any number of insertions of
the electron--external field vertex \eqref{eq:13.6.14} in the electron line. (See Figure
\ref{fig:weinberg-13-4}.) But as shown by the sum over permutations in Eq.~\eqref{eq:13.6.2}, these
diagrams in the external field approximation come from diagrams in the
underlying theory in which the photons attached to the electron line are
attached to the heavy charged particle line in all possible orders. (See
Figure \ref{fig:weinberg-13-5}.) The `uncrossed ladder' diagrams (labelled L) of Figure
\ref{fig:weinberg-13-5} do not dominate this sum unless the electron as well as the heavy
charged particle is non-relativistic. (These diagrams include contributions
from terms in old-fashioned perturbation theory whose intermediate states
contain the same particles as the initial and final states, leading to small
energy denominators when the electron and heavy charged particle are
both non-relativistic, while all other diagrams of Figure \ref{fig:weinberg-13-5} correspond
to intermediate states with either extra photons, electron--positron pairs,
or heavy particle--antiparticle pairs, which are suppressed by large energy
denominators.) The uncrossed ladders can be summed by solving an
integral equation, known as the Bethe--Salpeter equation~\hyperref[ref:13.10]{[10]}, but there is no
rationale for selecting out this subset of diagrams unless both particles are
non-relativistic, in which case the Bethe--Salpeter equation just reduces
to the ordinary non-relativistic Schr\"odinger equation, plus relativistic
corrections associated with the spin--orbit coupling that can be treated as
small perturbations. It must be said that the theory of relativistic effects
and radiative corrections in bound states is not yet in entirely satisfactory
shape.

In the derivation of the external field \eqref{eq:13.6.16} we evaluated the interaction of the heavy charged particle with the electromagnetic field only to
leading order in the photon momentum. There are corrections of higher
order in the photon momentum arising from the heavy particle's magnetic
dipole moment, electric quadrupole moment, etc. Also, of course, there
are radiative corrections arising from Feynman diagrams beyond those of
Figure \ref{fig:weinberg-13-4}, such as diagrams in which photons are emitted and absorbed
from the electron line or electron loops are inserted into photon lines. We
shall see in the next chapter that in bound states the diagrams of Figure
\ref{fig:weinberg-13-4} must be included to all orders, but all other corrections of higher
order in photon momenta or $e$ may be included as perturbations to these
diagrams.

% Figure 13.4: successive external-field insertions on an electron line.
% Source: physical PDF p. 584 (printed p. 561).
\begingroup
\renewcommand{\thefigure}{13.\arabic{figure}}
\setcounter{figure}{3}
\begin{figure}[tbp]
\centering
\begin{tikzpicture}[
  electron/.style={line width=0.7pt},
  field/.style={line width=0.62pt,decorate,
    decoration={snake,amplitude=0.95pt,segment length=4.7pt}},
  cross/.style={line width=0.62pt},
  every node/.style={font=\normalsize}
]
  \newcommand{\fieldinsertion}[2]{%
    \draw[field] (0,#1) -- (#2,#1);
    \draw[cross] (#2-0.09,#1-0.09) -- (#2+0.09,#1+0.09);
    \draw[cross] (#2-0.09,#1+0.09) -- (#2+0.09,#1-0.09);
  }
  \begin{scope}[xshift=-5.15cm]
    \draw[electron] (0,-1.45) -- (0,1.45);
    \fieldinsertion{0}{1.45}
  \end{scope}
  \node at (-2.82,0) {$+$};
  \begin{scope}[xshift=-1.45cm]
    \draw[electron] (0,-1.45) -- (0,1.45);
    \fieldinsertion{0.63}{1.45}
    \fieldinsertion{-0.63}{1.45}
  \end{scope}
  \node at (0.88,0) {$+$};
  \begin{scope}[xshift=2.25cm]
    \draw[electron] (0,-1.45) -- (0,1.45);
    \fieldinsertion{0.77}{1.45}
    \fieldinsertion{0}{1.45}
    \fieldinsertion{-0.77}{1.45}
  \end{scope}
  \node at (4.60,0) {$+\ \cdots$};
\end{tikzpicture}
\caption{Diagrams for the scattering of an electron by an external electromagnetic field. Here straight lines represent the electron; wavy lines ending in crosses represent its interaction with an external field.}
\label{fig:weinberg-13-4}
\end{figure}
\endgroup

% Figure 13.5: one-, two-, and three-photon exchanges with a heavy target.
% Source: physical PDF p. 584 (printed p. 561).
\begingroup
\renewcommand{\thefigure}{13.\arabic{figure}}
\setcounter{figure}{4}
\begin{figure}[tbp]
\centering
\begin{tikzpicture}[
  electron/.style={line width=0.7pt},
  target/.style={double,double distance=2.2pt,line width=0.52pt},
  photon/.style={line width=0.62pt,decorate,
    decoration={snake,amplitude=0.82pt,segment length=4.1pt}},
  every node/.style={font=\small}
]
  % Each panel spans x=0..2.15 and y=-1.12..1.12.
  \newcommand{\rails}{%
    \draw[electron] (0,-1.12) -- (0,1.12);
    \draw[target] (2.15,-1.12) -- (2.15,1.12);
  }

  % One-photon exchange (uncrossed ladder).
  \begin{scope}[xshift=-5.25cm,yshift=3.25cm]
    \rails
    \draw[photon] (0,0) -- (2.15,0);
    \node at (1.08,-1.68) {L};
  \end{scope}
  % Two-photon uncrossed ladder.
  \begin{scope}[xshift=-1.08cm,yshift=3.25cm]
    \rails
    \draw[photon] (0,0.55) -- (2.15,0.55);
    \draw[photon] (0,-0.55) -- (2.15,-0.55);
    \node at (1.08,-1.68) {L};
  \end{scope}
  % Two-photon crossed ladder.
  \begin{scope}[xshift=3.08cm,yshift=3.25cm]
    \rails
    \draw[photon] (0,0.55) -- (2.15,-0.55);
    \draw[photon] (0,-0.55) -- (2.15,0.55);
  \end{scope}

  % Three-photon uncrossed ladder.
  \begin{scope}[xshift=-5.25cm,yshift=0cm]
    \rails
    \draw[photon] (0,0.72) -- (2.15,0.72);
    \draw[photon] (0,0) -- (2.15,0);
    \draw[photon] (0,-0.72) -- (2.15,-0.72);
    \node at (1.08,-1.68) {L};
  \end{scope}
  % One upper rung, with the two lower rungs crossed.
  \begin{scope}[xshift=-1.08cm,yshift=0cm]
    \rails
    \draw[photon] (0,0.72) -- (2.15,0.72);
    \draw[photon] (0,0) -- (2.15,-0.72);
    \draw[photon] (0,-0.72) -- (2.15,0);
  \end{scope}
  % Two upper rungs crossed, with one lower straight rung.
  \begin{scope}[xshift=3.08cm,yshift=0cm]
    \rails
    \draw[photon] (0,0.72) -- (2.15,0);
    \draw[photon] (0,0) -- (2.15,0.72);
    \draw[photon] (0,-0.72) -- (2.15,-0.72);
  \end{scope}

  % The three remaining orderings.
  \begin{scope}[xshift=-5.25cm,yshift=-3.25cm]
    \rails
    \draw[photon] (0,0.72) -- (2.15,-0.72);
    \draw[photon] (0,0) -- (2.15,0.72);
    \draw[photon] (0,-0.72) -- (2.15,0);
  \end{scope}
  \begin{scope}[xshift=-1.08cm,yshift=-3.25cm]
    \rails
    \draw[photon] (0,0.72) -- (2.15,0);
    \draw[photon] (0,0) -- (2.15,-0.72);
    \draw[photon] (0,-0.72) -- (2.15,0.72);
  \end{scope}
  \begin{scope}[xshift=3.08cm,yshift=-3.25cm]
    \rails
    \draw[photon] (0,0.72) -- (2.15,-0.72);
    \draw[photon] (0,0) -- (2.15,0);
    \draw[photon] (0,-0.72) -- (2.15,0.72);
  \end{scope}
\end{tikzpicture}
\caption{Diagrams for the scattering of an electron by a heavy, charged target particle, which in the limit of large target mass yield the same result as the diagrams of Figure~\ref{fig:weinberg-13-4}. Here the single straight line is the electron; the double straight line is the heavy target particle; and wavy lines are virtual photons. Diagrams marked `L' are called uncrossed ladder graphs; they dominate the sum when the electron as well as the target particle is non-relativistic.}
\label{fig:weinberg-13-5}
\end{figure}
\endgroup

